\problem{Gaz parfait de fermions ultra-relativiste en trois dimension}

\question{Calculez l'énergie de Fermi en fonction de la densité n = N/V du gaz.}

L'énergie d'un fermion ultra-relativiste peut être réécrit comme 
\[
    \epsilon_p &= cp \notag \\
    &= c\hbar \mathbf{k}
\]

Pour l'énergie de Fermi, nous avons :
\[
    \varepsilon_F = c\hbar k_F
\]

, avec $k_F$ le vecteur d'onde de Fermi.

Or, nous savons que le vecteur d'onde de Fermi pour un gaz de fermion est
\[
    k_F = \qty(\frac{6\pi^2n}{g})^{1/3}
\]

, avec n la densité du gaz. En imbriquant les deux équations, nous avons la solution :

\solution{energie1}{\varepsilon_F = c\hbar \qty(\frac{6\pi^2n}{g})^{1/3}}

\question{ Calculez l'énergie moyenne du gaz en fonction de $\varepsilon_F$ et de N, à T = 0.}

Nous savons que 
\[
    E &= \sum_{\lambda}\epsilon_{\lambda}\bar{n}_{\lambda}^F \notag \\
    &= \sum_{k,s}\epsilon_{k,s}\bar{n}_{k,s}^F.
\]

Puisqu'il n'y a pas de champ magnétique, la somme sur s est indépendente du reste du calcul. 
Nous avons donc :
\[
    E &= \frac{gV}{8\pi^3}\int d^3k \epsilon_{k,s}\bar{n}_{k,s}^F \notag \\
    &= \frac{gV}{8\pi^3}\int d^3k c\hbar \mathbf{k} \Theta(k_F - \mathbf{k}).
\]
Réexprimons l'intégrale en coordonnées sphérique :
\[
    E &= \frac{4\pi V g}{8\pi^3} \int_{0}^{k_F} dk c\hbar k^3 \notag \\
    &= \frac{4\pi V g}{8\pi^3} \frac{c\hbar k_F^4}{4} \notag \\
    &= \frac{V}{8\pi^2}\varepsilon_F k_F^3 g \notag \\
    &= \frac{6\pi^2 N}{gV} \frac{Vg}{8\pi^2}\varepsilon_F \notag \\
    &= \frac{3}{4} N \varepsilon_F.
\]

Nous avons donc la solution :
\solution{energie2}{E = \frac{3}{4} N \varepsilon_F}

\question{Montrez que la pression (à une température quelconque) est donnée par p = E/3V, où E est l'énergie totale moyenne du gaz. Comparez ce résultat avec celui du gaz non-relativiste}

L'expression de la pression est :
\[
    p = -\sum_{\lambda} \pdv{\epsilon_{\lambda}}{V} \bar{n}_{\lambda}^F.
\]
Concentrons nous sur la dérivée. Nous obtenons :
\[
    \pdv{\epsilon_{\lambda}}{V} &= c\hbar\qty(\frac{6\pi^2N}{g})^{1/3} \pdv{V}\frac{1}{V^{1/3}} \notag \\
    &= c\hbar\qty(\frac{6\pi^2N}{Vg})^{1/3} \frac{-1}{3V} \notag \\
    &= \frac{-\epsilon_{\lambda}}{3V}. 
\]

En rempaçant cette équation dans l'équation initiale de la pression, nous avons :
\[
    p &= \frac{1}{3} \sum_{\lambda} \epsilon_{\lambda} \bar{n}_{\lambda}^F \notag \\
    &= \frac{E}{3V}.
\]

Nous avons donc une expression tel que :
\solution{sol1_3}{p = \frac{E}{3V}}

En comparaision avec avec un gaz de fermion non relativiste, nous remarquons que la différence est un facteur 2. 
Cela peux s'expliquer par l'expression de l'énergie qui n'est pas la même entre les deux gaz. Dans un gaz classique, 
l'énergie dépend en $k^2$ alors qu'elle dépend en k.

\question{Calculez la susceptibilité de Pauli à température nulle.}

Nous avons d'abord l'aimantation tel que dans les notes de cours :
\[
    M = \frac{\mu_Bg}{6\pi^2}\qty(k_{\downarrow}^3 - k_{\uparrow}^3).
\]

Nous savons que dans notre situation $k_{\uparrow} = \frac{\varepsilon_F - \mu_B B}{c\hbar}$ et $k_{\downarrow} = \frac{\varepsilon_F + \mu_B B}{c\hbar}$. 
Avec cela, nous pouvons continuer la résolution :
\[
    M &= \frac{\mu_Bg}{6\pi^2}\qty(\qty(\frac{\varepsilon_F + \mu_B B}{c\hbar})^3 - \qty(\frac{\varepsilon_F - \mu_B B}{c\hbar})^3) \notag \\
    &= \frac{\mu_Bg}{6\pi^2} \varepsilon^3 \qty(\qty(1 + \frac{\mu_B B}{\varepsilon_Fc\hbar})^3 - \qty(1 - \frac{\mu_B B}{\varepsilon_Fc\hbar})^3)
\]

Puisque nous sommes dans un solide, nous pouvons approximer $\mu_BB << \varepsilon_F$. 
Nous pouvons alors développer l'équation de l'aimantation comme
\[
    M &= \frac{\mu_Bg}{6\pi^2} \varepsilon^3 \qty(1 + \frac{3\mu_B B}{\varepsilon_Fc\hbar} - 1 + \frac{3\mu_B B}{\varepsilon_Fc\hbar}) \notag \\
    &= \frac{\mu_Bg}{6\pi^2} \varepsilon^3 \frac{6\mu_B B}{\varepsilon_Fc\hbar} \notag \\
    &= \frac{6\mu_B^2B}{\varepsilon_F}n.
\]

Le lien entre la susceptibilité magnétique et l'aimantation est :
\[
    M = \chi B.
\]

Nous avons donc la réponse :
\[
    \chi = \frac{M}{B} \notag
\]
\solution{1.4}{\chi = \frac{6\mu_B^2B}{\varepsilon_F}n}
\clearpage

